%==================================================================
% ABSTRAK (Bahasa Indonesia)
% Acuan: Pedoman Lampiran 9A
%   - judul TA 14 pt tebal, rata tengah, spasi 1
%   - nama, NIM, pembimbing 2 spasi di bawah judul
%   - kata ABSTRAK 2 spasi di bawah nama, 14 pt tebal
%   - isi abstrak spasi 1, baris pertama masuk 1,25 cm
%   - maksimal 350 kata, kata kunci 3-5 kata
%==================================================================
\clearpage
\thispagestyle{plain}
\frontmatterentry{ABSTRAK}

\begin{spacing}{1}
\begin{center}
{\fontsize{14}{16.8}\selectfont\bfseries \MakeUppercase{\judulid}\par}
\end{center}

\vspace{2\baselineskip}

\noindent
\begin{tabular}{@{}l@{\hspace{0.6em}}c@{\hspace{0.6em}}p{10.8cm}@{}}
Nama mahasiswa & : & \penulis \\
NIM            & : & \nim \\
Pembimbing     & : & 1.\ \pembimbingA \\
               &   & 2.\ \pembimbingB \\
\end{tabular}

\vspace{2\baselineskip}

\begin{center}
{\fontsize{14}{16.8}\selectfont\bfseries ABSTRAK\par}
\end{center}

\vspace{2\baselineskip}

%--- Isi abstrak: ganti dengan ringkasan penelitian Anda -------------
Penurunan mutu kesegaran ikan terjadi sangat cepat setelah ikan
ditangkap, sehingga penilaian kesegaran secara cepat dan objektif
menjadi kebutuhan penting bagi pembeli maupun pelaku usaha perikanan.
Selama ini penilaian kesegaran masih dilakukan secara organoleptik
oleh tenaga terlatih sehingga hasilnya bersifat subjektif dan sulit
dilakukan dalam jumlah besar. Penelitian ini bertujuan merancang
bangun sistem deteksi kesegaran ikan berbasis pengolahan citra digital
dan jaringan saraf tiruan yang mampu mengklasifikasikan tingkat
kesegaran ikan secara otomatis. Citra mata dan insang ikan diambil
menggunakan kamera digital pada kotak pengambilan gambar dengan
pencahayaan tetap, kemudian diolah melalui tahap penyekalaan ukuran,
perbaikan kontras, dan ekstraksi ciri warna pada ruang warna HSV.
Ciri yang diperoleh digunakan sebagai masukan jaringan saraf tiruan
dengan arsitektur perambatan balik dua lapisan tersembunyi. Pengujian
dilakukan terhadap 300 citra ikan yang terdiri atas tiga kelas
kesegaran, yaitu segar, cukup segar, dan tidak segar. Hasil pengujian
menunjukkan akurasi klasifikasi sebesar 93,33\% dengan waktu
pemrosesan rata-rata 0,42 detik per citra. Nilai presisi dan perolehan
tertinggi diperoleh pada kelas ikan segar, sedangkan kesalahan
terbanyak terjadi antara kelas cukup segar dan tidak segar karena
kemiripan warna insang. Dengan demikian sistem yang dirancang dapat
digunakan sebagai alternatif penilaian kesegaran ikan secara cepat,
objektif, dan tidak merusak.

\vspace{3\baselineskip}

\noindent
\textbf{Kata kunci:} jaringan saraf tiruan, kesegaran ikan, pengolahan
citra digital, ruang warna HSV
%-------------------------------------------------------------------
\end{spacing}
